\section*{Bab \ref{ch:b1:findim} --- Dimensi Hingga}

\begin{solution}{exo:b1:findim:1}
\begin{enumerate}
  \item $z = -x - y$: $F = \{(x, y, -x-y)\} =
        \operatorname{Vect}\bigl((1,0,-1), (0,1,-1)\bigr)$; kedua
        vektornya \omterm{def:b1:vspaces:free}{bebas} (\omterm{prop:b1:vspaces:coordinates}{koordinat}): $\dim F = 2$.
  \item $G = \{(x, x, 2t, t)\} =
        \operatorname{Vect}\bigl((1,1,0,0), (0,0,2,1)\bigr)$: \omterm{def:b1:vspaces:free}{bebas},
        $\dim G = 2$.
  \item $P(1) = P'(1) = 0$ berarti $(X-1)^2 \mid P$
        (\cref{prop:b1:poly:multiplicity}): $P = (X-1)^2(aX + b)$.
        \omterm{def:b1:vspaces:free}{Basis} $\bigl((X-1)^2, X(X-1)^2\bigr)$, berdimensi $2$.
\end{enumerate}
\end{solution}

\begin{solution}{exo:b1:findim:2}
$(3,5,7) = (1,1,1) + 2(1,2,3)$ dan $(0,1,2) = (1,2,3) - (1,1,1)$:
keduanya kombinasi dua yang pertama, yang \omterm{def:b1:vspaces:free}{bebas} (karena tak
sebanding). \omterm{def:b1:findim:rank}{Rank} $2$; \omterm{def:b1:vspaces:free}{basis} \omterm{def:b1:vspaces:span}{rentangnya}:
$\bigl((1,1,1), (1,2,3)\bigr)$.
\end{solution}

\begin{solution}{exo:b1:findim:3}
Cobalah menambahkan $e_1 = (1,0,0,0)$ dan $e_3 = (0,0,1,0)$. Keluarga
$\bigl((1,1,0,0), (0,0,1,1), e_1, e_3\bigr)$ bersifat \omterm{def:b1:vspaces:free}{bebas}: karena kombinasi
nol $\alpha(1,1,0,0) + \beta(0,0,1,1) + \gamma e_1 + \delta
e_3 = 0$ terbaca $(\alpha + \gamma, \alpha, \beta + \delta, \beta) =
0$, sehingga $\alpha = \beta = 0$, lalu $\gamma = \delta = 0$. Empat vektor
yang \omterm{def:b1:vspaces:free}{bebas} dalam dimensi $4$: yaitu sebuah \omterm{def:b1:vspaces:free}{basis}
(\cref{prop:b1:findim:twoofthree}).
\end{solution}

\begin{solution}{exo:b1:findim:4}
Derajat $0, 1, 2, 3$ berbeda berpasangan: jadi \omterm{def:b1:vspaces:free}{bebas}
(\cref{prop:b1:vspaces:freecriteria}), empat vektor dalam dimensi
$4$: jadi \omterm{def:b1:vspaces:free}{basis}. Untuk $X^3$: jabarkanlah ke bawah,
\[
X(X-1)(X-2) = X^3 - 3X^2 + 2X,
\qquad
X(X-1) = X^2 - X,
\]
sehingga $X^3 - X(X-1)(X-2) = 3X^2 - 2X$; dan $3X^2 - 2X = 3(X^2 - X) + X
= 3\,X(X-1) + X$. Maka
\[
X^3 = X(X-1)(X-2) + 3\,X(X-1) + 1\cdot X + 0\cdot 1 :
\]
\omterm{prop:b1:vspaces:coordinates}{koordinat} $(0, 1, 3, 1)$ pada $\bigl(1, X, X(X-1),
X(X-1)(X-2)\bigr)$. (Inilah bilangan Stirling dalam samaran.)
\end{solution}

\begin{solution}{exo:b1:findim:5}
Grassmann: $\dim(F \cap G) = 4 + 5 - \dim(F + G)$, dan $F + G$ merupakan
\omterm{def:b1:vspaces:subspace}{subruang} $E$ yang memuat $G$: $5 \leq \dim(F+G) \leq 7$. Maka
$\dim(F \cap G) \in \{2, 3, 4\}$. Perwujudannya di dalam $\R^7$ dengan
\omterm{def:b1:vspaces:free}{basis} kanonik $(e_1, \dots, e_7)$, mengambil $G =
\operatorname{Vect}(e_1, \dots, e_5)$:
\begin{itemize}
  \item $F = \operatorname{Vect}(e_1, e_2, e_6, e_7)$: $F + G = \R^7$,
        irisannya $\operatorname{Vect}(e_1, e_2)$, berdimensi $2$;
  \item $F = \operatorname{Vect}(e_1, e_2, e_3, e_6)$: irisan
        yang berdimensi $3$;
  \item $F = \operatorname{Vect}(e_1, e_2, e_3, e_4) \subseteq G$:
        berdimensi $4$.
\end{itemize}
\end{solution}

\begin{solution}{exo:b1:findim:6}
$H$ merupakan \omterm{def:b1:vspaces:subspace}{subruang} (\cref{exo:b1:vspaces:1}~(4) secara harfiah). Menurut
teorema faktor (\cref{thm:b1:poly:factor}), $P \in H \iff P =
(X-1)Q$ dengan $\deg Q \leq n - 1$: jadi \omterm{def:b1:logic:map}{pemetaan} $Q \mapsto (X-1)Q$ merupakan
bijeksi linear dari $\R_{n-1}[X]$ pada $H$, sehingga \omterm{def:b1:vspaces:free}{basis} $H$ adalah
\[
\bigl((X-1),\ (X-1)X,\ (X-1)X^2,\ \dots,\ (X-1)X^{n-1}\bigr),
\qquad \dim H = n .
\]
Sebuah garis \omterm{def:b1:vspaces:sum}{pelengkap}: $\operatorname{Vect}(1)$ (yaitu konstanta). Memang
$H \cap \operatorname{Vect}(1) = \{0\}$ (karena konstanta taknol tak
lenyap di $1$) dan dimensinya berjumlah $n + 1$: jadi menurut Grassmann, $H
\oplus \operatorname{Vect}(1) = \R_n[X]$.
\end{solution}

\begin{solution}{exo:b1:findim:7}
Penghapusan: dengan membuang $u_p$, \omterm{def:b1:vspaces:span}{rentangnya} hanya dapat menyusut; dan ia menyusut
paling banyak satu dimensi, karena menambahkan $u_p$ kembali pada \omterm{def:b1:vspaces:free}{basis}
\omterm{def:b1:vspaces:span}{rentang} yang lebih kecil memberikan keluarga pembangun bagi yang lebih besar dengan paling
banyak satu vektor tambahan. Secara setangkup, menambahkan sebuah vektor $v$: \omterm{def:b1:vspaces:span}{rentang}
yang baru memuat yang lama dengan paling banyak satu pembangun tambahan, sehingga
dimensinya adalah $r$ (jika $v$ sudah ada di dalam \omterm{def:b1:vspaces:span}{rentangnya}) atau $r + 1$
(jika tidak, karena menurut \cref{prop:b1:vspaces:freecriteria}~(2) sebuah \omterm{def:b1:vspaces:free}{basis}
dapat diperluas). Kedua \omterm{def:b1:logic:statement}{pernyataan} itu bersama-sama memberikan kesimpulan ``\omterm{def:b1:findim:rank}{rank}
bersifat $1$-Lipschitz''.
\end{solution}

\begin{solution}{exo:b1:findim:8}
Sepanjang rantai yang naik tegas, dimensinya naik tegas
($F_i \subseteq F_{i+1}$, $F_i \neq F_{i+1}$ dan
\cref{thm:b1:findim:subspaces}: karena kesamaan dimensinya akan memaksa
kesamaan ruangnya). Jadi $\dim F_1 < \dim F_2 < \dots < \dim F_k$ merupakan
barisan bilangan bulat yang naik tegas di dalam $\intint{0}{n}$: yaitu paling
banyak $n + 1$ nilai, sehingga $k \leq n + 1$. Rantai maksimal di dalam $\R^n$:
\[
\{0\} \subsetneq \operatorname{Vect}(e_1) \subsetneq
\operatorname{Vect}(e_1, e_2) \subsetneq \dots \subsetneq \R^n ,
\]
yang berpanjang tepat $n + 1$.
\end{solution}

\begin{solution}{exo:b1:findim:9}
$F + G$ memuat $F$ secara tegas (karena $G \not\subseteq F$), sehingga
$\dim(F + G) = n$ dan Grassmann memberikan $\dim(F \cap G) = (n-1) +
(n-1) - n = n - 2$.

Klaim umumnya, lewat induksi pada $k$: bahwa irisan $I_k$ atas $k$
hiperbidang memenuhi $\dim I_k \geq n - k$. Benar untuk $k = 1$. Langkahnya: $I_{k+1}
= I_k \cap H_{k+1}$, dan Grassmann di dalam $E$:
\[
\dim(I_k \cap H_{k+1}) = \dim I_k + (n - 1) - \dim(I_k + H_{k+1})
\geq \dim I_k + (n-1) - n \geq n - k - 1 . \qedhere
\]
\end{solution}

\begin{solution}{exo:b1:findim:10}
\begin{enumerate}
  \item Syarat $u_{n+2} - 3u_{n+1} + 2u_n = 0$ bersifat linear
        dan dipenuhi barisan nol: jadi $E$ merupakan \omterm{def:b1:vspaces:subspace}{subruang}. Lewat
        induksi, $u_0$ dan $u_1$ menentukan setiap $u_n$, dan
        kebergantungannya linear (karena setiap langkahnya kombinasi linear atas
        kedua nilai sebelumnya); sebaliknya setiap pasangan $(a, b)$
        muncul dari tepat satu barisan $E$ (definisikanlah $u_n$ lewat
        rekurensinya). Seperti pada \cref{ex:b1:findim:dims}, $E$
        terparameterkan secara \omterm{def:b1:logic:inj}{bijektif} dan linear oleh $(u_0, u_1) \in
        \R^2$: sehingga $\dim E = 2$.
  \item Konstanta: $3\cdot1 - 2\cdot1 = 1$. Geometri: $3\cdot
        2^{n+1} - 2\cdot 2^n = (6 - 2)2^n = 2^{n+2}$. Keduanya ada di dalam
        $E$. Kebebasannya: $a\cdot 1 + b\cdot 2^n = 0$ untuk semua $n$
        memberikan, di $n = 0$ dan $n = 1$: $a + b = 0$, $a + 2b = 0$,
        jadi $a = b = 0$. Dua vektor \omterm{def:b1:vspaces:free}{bebas} dalam dimensi $2$: yaitu \omterm{def:b1:vspaces:free}{basis}
        (\cref{prop:b1:findim:twoofthree}).
  \item Pecahkanlah $a + b = 0$, $a + 2b = 1$: $b = 1$, $a = -1$, sehingga
        $u_n = 2^n - 1$ (yaitu barisan Mersenne).
\end{enumerate}
\end{solution}

\begin{solution}{exo:b1:findim:11}
\begin{enumerate}
  \item Grassmann: $\dim(F \cap G) = \dim F + \dim G - \dim(F +
        G) \geq \dim F + \dim G - n > 0$, sehingga $F \cap G \neq
        \{0\}$. Dengan $n = 100$: $51 + 50 - 100 = 1 > 0$, jadi
        irisannya memuat sebuah garis.
  \item Jika $F \cap H = \{0\}$, maka jumlahnya langsung dan $\dim F +
        (n - 1) = \dim(F \oplus H) \leq n$, sehingga $\dim F \leq 1$.
        (Sebaliknya sebuah garis yang tak termuat di dalam $H$ memang
        memenuhi ini: karena hiperbidang nyaris tak melewatkan apa pun.)
\end{enumerate}
\end{solution}

\begin{solution}{exo:b1:findim:12}
Induksi menurun pada $d = \dim F = \dim G$, dari $d = n$ ke $d =
0$. Jika $d = n$: maka $F = G = E$ dan $S = \{0\}$ berhasil. Andaikan
\omterm{def:b1:logic:statement}{pernyataannya} berlaku bagi pasangan \omterm{def:b1:vspaces:subspace}{subruang} berdimensi $d + 1 \leq n$,
lalu misalkan $\dim F = \dim G = d < n$.

Jika $F = G$: ambillah sebagai $S$ sebarang \omterm{def:b1:vspaces:sum}{pelengkap} $F$
(\cref{thm:b1:findim:subspaces}). Jika $F \neq G$: maka keduanya sejati,
sehingga menurut \cref{exo:b1:vspaces:12} ada $x \notin F \cup G$.
Jumlah $F' = F \oplus \operatorname{Vect}(x)$ dan $G' = G \oplus
\operatorname{Vect}(x)$ bersifat langsung (karena $x \notin F$, $x \notin G$)
dan berdimensi $d + 1$; jadi menurut hipotesis induksinya keduanya
mempunyai \omterm{def:b1:vspaces:sum}{pelengkap} bersama $S'$: $E = F' \oplus S' = G' \oplus
S'$. Tetapkanlah $S = \operatorname{Vect}(x) \oplus S'$ --- yang langsung,
karena $S' \cap \operatorname{Vect}(x) \subseteq S' \cap F' =
\{0\}$.

Maka $F + S = F + \operatorname{Vect}(x) + S' = F' + S' = E$, dan
$F \cap S = \{0\}$: karena jika $f = \lambda x + s'$ dengan $f \in F$, $s'
\in S'$, maka $s' = f - \lambda x \in F' \cap S' = \{0\}$, sehingga $f
= \lambda x$, yang memaksa $\lambda = 0$ (karena $x \notin F$) dan $f = 0$.
Jadi $E = F \oplus S$, dan secara setangkup $E = G \oplus S$.
\end{solution}

\begin{solution}{pb:b1:findim:1}
\textbf{1.} $\Q$ memuat $0 \neq 1$, stabil terhadap penjumlahan,
perkalian dan lawan, dan setiap rasional taknol mempunyai
invers rasional: jadi sebuah \omterm{def:b1:structures:field}{lapangan} (\cref{def:b1:structures:field}). Adapun
definisi \omterm{def:b1:vspaces:span}{rentang}, kebebasan, \omterm{def:b1:vspaces:free}{basis} dan bukti Bab
18--19 hanya memakai penjumlahan vektor, kedistributifan dan
aritmetika skalar di dalam sebuah \omterm{def:b1:structures:field}{lapangan} --- tak pernah nilai mutlak, urutan atau
limit. Jadi $\R$, dengan penjumlahannya sendiri dan perkalian
$\Q \times \R \to \R$, merupakan \omterm{def:b1:vspaces:def}{ruang vektor} atas $\Q$, dan semua
teoremanya berlaku. Pembagian dengan sebuah skalar terjadi sekali pada bukti
\cref{thm:b1:findim:exchange}: yaitu untuk ``menyelesaikan $g_k$'' orang \omterm{def:b1:arith:divides}{membagi}
dengan koefisiennya yang taknol.

\textbf{2.} Jika $a + b\sqrt2 = 0$ dengan $a, b \in \Q$ yang tak keduanya
nol: maka $b \neq 0$ akan memberikan $\sqrt2 = -a/b \in \Q$, yang bertentangan dengan
keirasionalan $\sqrt2$ (\cref{exo:b1:arith:7} dengan $p =
2$); sehingga $b = 0$, lalu $a = 0$: jadi \omterm{def:b1:vspaces:free}{bebas} atas $\Q$. Sedangkan atas $\R$
hubungan $\sqrt2\cdot 1 + (-1)\cdot\sqrt2 = 0$ tak sepele:
jadi terkait. Maka $\dim_\Q \Q(\sqrt2) = 2$, dengan \omterm{prop:b1:vspaces:coordinates}{koordinat} tunggal
$(a, b)$.

\textbf{3.} $(a + b\sqrt2)(c + d\sqrt2) = (ac + 2bd) + (ad +
bc)\sqrt2 \in \Q(\sqrt2)$. Maka
\[
\sigma(uv) = (ac + 2bd) - (ad + bc)\sqrt2
= (a - b\sqrt2)(c - d\sqrt2) = \sigma(u)\sigma(v),
\]
sehingga $N(uv) = uv\,\sigma(uv) = u\sigma(u)\,v\sigma(v) = N(u)N(v)$.
Jika $N(u) = a^2 - 2b^2 = 0$ dengan $u \neq 0$: maka $b \neq 0$ akan memberikan
$(a/b)^2 = 2$, yaitu akar kuadrat rasional atas $2$; sehingga $b = 0$, lalu $a
= 0$ dan $u = 0$: bertentangan. Maka $N(u) \neq 0$ untuk $u \neq
0$.

\textbf{4.} $u \cdot \dfrac{\sigma(u)}{N(u)} = \dfrac{N(u)}{N(u)}
= 1$, dan $\sigma(u)/N(u) \in \Q(\sqrt2)$: jadi setiap unsur taknol
dapat dibalik di dalam $\Q(\sqrt2)$, yang karenanya merupakan sublapangan
$\R$. Contohnya: $N(3 + 2\sqrt2) = 9 - 8 = 1$, sehingga
\[
\frac1{3 + 2\sqrt2} = 3 - 2\sqrt2 ;
\qquad
N(1 + \sqrt2) = -1, \quad
\frac1{1 + \sqrt2} = \sqrt2 - 1 .
\]

\textbf{5.} Jika $\sqrt6 = p/q$ maka $6q^2 = p^2$; sehingga pangkat
$2$ adalah $1 + 2v_2(q)$, yang ganjil, di kiri, dan $2v_2(p)$, yang genap, di
kanan: mustahil. Kini andaikan $\sqrt3 = a + b\sqrt2$ dengan
$a, b \in \Q$. Mengkuadratkannya: $3 = a^2 + 2b^2 + 2ab\sqrt2$. Jika $ab
\neq 0$, maka $\sqrt2 = (3 - a^2 - 2b^2)/(2ab) \in \Q$:
mustahil. Jika $b = 0$: maka $\sqrt3 = a \in \Q$, yang bertentangan dengan
\cref{exo:b1:arith:7} ($p = 3$). Jika $a = 0$: maka $\sqrt3 = b\sqrt2$,
dan dengan mengalikannya dengan $\sqrt2$: $\sqrt6 = 2b \in \Q$: mustahil.
Jadi $\sqrt3 \notin \Q(\sqrt2)$.

\textbf{6.} Hasil kali vektor basisnya adalah
\[
\sqrt2\,\sqrt3 = \sqrt6,\quad
\sqrt2\,\sqrt6 = 2\sqrt3,\quad
\sqrt3\,\sqrt6 = 3\sqrt2,\quad
(\sqrt2)^2 = 2,\ (\sqrt3)^2 = 3,\ (\sqrt6)^2 = 6,
\]
semuanya di dalam $M$. Adapun hasil kali dua unsur $M$ terjabarkan lewat
kebilinearan menjadi kombinasi rasional atas ini: jadi $M$ stabil
terhadap perkalian.

\textbf{7.} Misalkan $x + y\sqrt3 = 0$ dengan $x, y \in K = \Q(\sqrt2)$.
Jika $y \neq 0$, maka $\sqrt3 = -x/y \in K$ (menurut pertanyaan 4: $K$ sebuah
\omterm{def:b1:structures:field}{lapangan}), yang bertentangan dengan pertanyaan 5. Jadi $y = 0$, lalu $x = 0$: sehingga $(1,
\sqrt3)$ \omterm{def:b1:vspaces:free}{bebas} atas $K$. Ia membangun: $K + K\sqrt3 =
\{(a + b\sqrt2) + (c + d\sqrt2)\sqrt3\} = \operatorname{Vect}_\Q
(1, \sqrt2, \sqrt3, \sqrt6) = M$. Maka $\dim_K M = 2$.

\textbf{8.} Sebuah hubungan $a + b\sqrt2 + c\sqrt3 + d\sqrt6 = 0$
terkelompokkan ulang menjadi $(a + b\sqrt2) + (c + d\sqrt2)\sqrt3 = 0$ dengan
koefisien di dalam $K$; jadi menurut pertanyaan 7 keduanya lenyap, dan menurut pertanyaan 2
$a = b = 0$ dan $c = d = 0$. Sehingga keluarganya \omterm{def:b1:vspaces:free}{bebas} dan $\dim_\Q M
= 4$.

\textbf{9.} Setiap $x \in M$ tertulis $x = \sum_j y_j f_j$ dengan $y_j
\in K$ (yaitu \omterm{def:b1:vspaces:free}{basis} $M$ atas $K$), dan setiap $y_j = \sum_i
\lambda_{ij} e_i$ dengan $\lambda_{ij} \in \Q$ (yaitu \omterm{def:b1:vspaces:free}{basis} $K$ atas
$\Q$); dengan menyulihkannya, $x = \sum_{i,j} \lambda_{ij}\, e_i f_j$: jadi
hasil kalinya membangun $M$ atas $\Q$.

\textbf{10.} Andaikan $\sum_{i,j} \lambda_{ij}\, e_i f_j = 0$ dengan
$\lambda_{ij} \in \Q$. Kelompokkanlah ulang: $\sum_j \bigl(\sum_i \lambda_{ij}
e_i\bigr) f_j = 0$, dan jumlah dalamnya termasuk $K$. Kebebasan
$(f_j)$ atas $K$ memberikan $\sum_i \lambda_{ij} e_i = 0$ untuk setiap
$j$; lalu kebebasan $(e_i)$ atas $\Q$ memberikan $\lambda_{ij} = 0$
untuk semua $i, j$.

\textbf{11.} Menurut pertanyaan 9 dan 10, $(e_i f_j)_{i \leq m,\, j
\leq n}$ merupakan \omterm{def:b1:vspaces:free}{basis} $M$ atas $\Q$ dengan $mn$ unsur:
\[
\dim_\Q M = m\,n = \dim_\Q K \cdot \dim_K M .
\]
Periksa: $\dim_\Q K = 2$ dan $\dim_K M = 2$ (pertanyaan 2 dan 7)
memberikan $\dim_\Q M = 4$, yaitu pertanyaan 8.

\textbf{12.} Vektor $u a_1, \dots, u a_d$ termasuk $A$
(karena kestabilannya). Semuanya \omterm{def:b1:vspaces:free}{bebas}: karena jika $\sum_i \lambda_i\, u a_i = 0$,
maka $u \sum_i \lambda_i a_i = 0$ di dalam $\R$, dan $u \neq 0$ memaksa
$\sum_i \lambda_i a_i = 0$, sehingga $\lambda_i = 0$ (karena $a_i$ membentuk
sebuah \omterm{def:b1:vspaces:free}{basis}). Jadi $(u a_1, \dots, u a_d)$ merupakan \omterm{def:b1:vspaces:free}{keluarga bebas} berisi $d$
vektor di dalam $A$, dengan $\dim_\Q A = d$: yaitu sebuah \omterm{def:b1:vspaces:free}{basis}
(\cref{prop:b1:findim:twoofthree}). Khususnya $1 \in A$
terurai menjadi $1 = \sum_i \mu_i\, u a_i = u\,v$ dengan $v = \sum_i
\mu_i a_i \in A$: sehingga invers $u$ ada di dalam $A$. Ini memulihkan
pertanyaan 4 ($A = \Q(\sqrt2)$) dan sekaligus membuktikan bahwa $M$
merupakan sublapangan $\R$ (bersama pertanyaan 6).

\textbf{13.} Kesemua $d + 1$ vektor $1, x, x^2, \dots, x^{d}$
ada di dalam $A$ (karena stabil terhadap hasil kali); sedangkan \omterm{def:b1:vspaces:free}{keluarga bebas} $A$ berisi
paling banyak $d$ vektor (\cref{prop:b1:findim:twoofthree}), sehingga semuanya
terkait: jadi ada rasional $\lambda_0, \dots, \lambda_d$,
yang tak semuanya nol, dengan $\sum_k \lambda_k x^k = 0$. Maka \omterm{def:b1:poly:def}{polinomial} $P
= \sum_k \lambda_k X^k$ taknol, berkoefisien rasional,
berderajat $\leq d$, dan $P(x) = 0$.

\textbf{14.} $s^2 = 2 + 2\sqrt6 + 3 = 5 + 2\sqrt6$: berkoordinat
$(5, 0, 0, 2)$. Lalu
\[
s^3 = s\,s^2 = (\sqrt2 + \sqrt3)(5 + 2\sqrt6)
= 5\sqrt2 + 2\sqrt{12} + 5\sqrt3 + 2\sqrt{18}
= 11\sqrt2 + 9\sqrt3 ,
\]
berkoordinat $(0, 11, 9, 0)$ (dengan memakai $\sqrt{12} = 2\sqrt3$,
$\sqrt{18} = 3\sqrt2$). Akhirnya $s^4 = (s^2)^2 = 25 + 20\sqrt6 +
24 = 49 + 20\sqrt6$: berkoordinat $(49, 0, 0, 20)$.

\textbf{15.} Andaikan $a + bs + cs^2 + ds^3 = 0$. Dengan membaca keempat
\omterm{prop:b1:vspaces:coordinates}{koordinatnya} pada $(1, \sqrt2, \sqrt3, \sqrt6)$:
\[
a + 5c = 0,\qquad b + 11d = 0,\qquad b + 9d = 0,\qquad 2c = 0 .
\]
Jadi $c = 0$, lalu $a = 0$; dengan mengurangkan persamaan tengahnya, $2d =
0$, lalu $b = 0$: sehingga $(1, s, s^2, s^3)$ \omterm{def:b1:vspaces:free}{bebas}. Empat vektor \omterm{def:b1:vspaces:free}{bebas}
di dalam $M$ yang berdimensi $4$: jadi sebuah \omterm{def:b1:vspaces:free}{basis}, sehingga
$\operatorname{Vect}_\Q(1, s, s^2, s^3) = M$. Dari pertanyaan 14,
$s^3 - 9s = 2\sqrt2$ dan $11s - s^3 = 2\sqrt3$:
\[
\sqrt2 = \frac{s^3 - 9s}{2},
\qquad
\sqrt3 = \frac{11s - s^3}{2} .
\]

\textbf{16.} $s^4 - 10s^2 + 1 = (49 + 20\sqrt6) - 10(5 +
2\sqrt6) + 1 = 0$. Adapun \omterm{def:b1:poly:def}{polinomial monik} berderajat $\leq 3$
yang lenyap di $s$ akan menghasilkan kombinasi nol yang tak sepele atas
$(1, s, s^2, s^3)$, yang bertentangan dengan pertanyaan 15: jadi $X^4 - 10X^2 + 1$
berderajat terkecil. Akarnya: $X^2 = 5 \pm 2\sqrt6 = (\sqrt3 \pm
\sqrt2)^2$, sehingga keempat akar realnya adalah $\pm(\sqrt3 + \sqrt2)$ dan
$\pm(\sqrt3 - \sqrt2)$, yaitu $\pm\sqrt2 \pm \sqrt3$.

\textbf{17.} Dari $s^4 - 10s^2 + 1 = 0$: $s\,(10s - s^3) = 1$, sehingga
\[
\frac1s = 10s - s^3 = 10(\sqrt2 + \sqrt3) - (11\sqrt2 + 9\sqrt3)
= \sqrt3 - \sqrt2 ,
\]
yang selaras dengan $(\sqrt3 + \sqrt2)(\sqrt3 - \sqrt2) = 1$. Jika $a +
b\sqrt2 + c\sqrt3 + d\sqrt6 = r \in \Q$, maka $(a - r) + b\sqrt2
+ c\sqrt3 + d\sqrt6 = 0$, dan kebebasannya (pertanyaan 8) memaksa $b = c
= d = 0$ (dan $a = r$). Adapun bagi $\sqrt2 + \sqrt3 + \sqrt6$
\omterm{prop:b1:vspaces:coordinates}{koordinatnya} $(0, 1, 1, 1)$ tak berbentuk demikian: jadi irasional ---
yaitu \omterm{def:b1:logic:statement}{pernyataan} yang lebih kuat daripada \cref{exo:b1:reals:7}, yang diperoleh tanpa
kiat pengkuadratan apa pun.

\textbf{18.} Sebuah akar rasional $p/q$ (dalam suku terkecil) atas $X^3 - 2$
memenuhi $p^3 = 2q^3$, dan kriteria
\cref{exo:b1:poly:5} memberikan $p \mid 2$, $q \mid 1$: jadi calonnya
$\pm1, \pm2$, yang pangkat tiganya $\pm1, \pm8 \neq 2$. Jadi tak ada akar
rasional; khususnya $t = 2^{1/3} \notin \Q$, sehingga $(1, t)$ \omterm{def:b1:vspaces:free}{bebas}
atas $\Q$ (seperti pada pertanyaan 2).

\textbf{19.} Andaikan $(1, t, t^2)$ terkait: maka ada $P \in
\Q[X]$ taknol dengan $\deg P \leq 2$ yang $P(t) = 0$; pilihlah $P$ yang demikian
berderajat \emph{terkecil} $d \geq 1$. Menurut pertanyaan 18, $d \neq 1$, sehingga $d
= 2$. Pembagian Euclid (\cref{thm:b1:poly:division}): $X^3 - 2
= PQ + R$ dengan $Q, R \in \Q[X]$, $\deg R < 2$. Menilainya di $t$:
$0 = P(t)Q(t) + R(t) = R(t)$, sehingga $R$ lenyap di $t$; lalu keminimalan
$d$ memaksa $R = 0$. Maka $X^3 - 2 = PQ$ dengan $\deg Q = 1$: sehingga
akar rasional $Q$ merupakan akar rasional $X^3 - 2$,
yang bertentangan dengan pertanyaan 18. Jadi $(1, t, t^2)$ \omterm{def:b1:vspaces:free}{bebas} dan
$\dim_\Q A = 3$. Kestabilannya: $t^3 = 2$ mereduksi setiap hasil kali
$1, t, t^2$ menjadi kombinasi atasnya ($t\cdot t^2 = 2$, $t^2\cdot
t^2 = 2t$); dan $A$ memuat $1$: jadi menurut pertanyaan 12, $A$ merupakan sublapangan
$\R$.

\textbf{20.} Andaikan $t \in M$. Maka $t^2 \in M$ (karena kestabilannya), sehingga
$A \subseteq M$, dan $M$ merupakan \omterm{def:b1:vspaces:def}{ruang vektor} atas \omterm{def:b1:structures:field}{lapangan} $A$:
karena aksiomanya adalah aksioma aritmetika $\R$, yang dibatasi. Ia
\omterm{def:b1:findim:def}{berdimensi hingga} atas $A$ (karena keluarga pembangun $\Q$ yang hingga
membangun a fortiori atas $A \supseteq \Q$). Adapun \omterm{pb:b1:findim:1}{hukum menara} bagi
$\Q \subseteq A \subseteq M$ memberikan
\[
4 = \dim_\Q M = \dim_\Q A \cdot \dim_A M = 3\,\dim_A M ,
\]
mustahil: karena $3$ tak \omterm{def:b1:arith:divides}{membagi} $4$. Jadi $2^{1/3} \notin M$: tak ada
kombinasi rasional atas $1, \sqrt2, \sqrt3, \sqrt6$ yang sama dengan
$2^{1/3}$.

\textbf{21.} $\Q(\sqrt2) \subseteq M$, sehingga $t \notin \Q(\sqrt2)$.
Jika $t + \sqrt2 = r \in \Q$, maka $t = r - \sqrt2 \in \Q(\sqrt2)$:
bertentangan --- jadi $2^{1/3} + \sqrt2$ irasional. Jika $t = a +
b\sqrt3$, maka $t \in \operatorname{Vect}_\Q(1, \sqrt3) \subseteq
M$: bertentangan lagi.

\textbf{22.} $B$ merupakan \omterm{def:b1:vspaces:def}{ruang vektor} atas \omterm{def:b1:structures:field}{lapangan} $A$
(yaitu pembatasan skalarnya), yang \omterm{def:b1:findim:def}{berdimensi hingga} karena $\dim_\Q B$
hingga. Adapun \omterm{pb:b1:findim:1}{hukum menara} bagi $\Q \subseteq A \subseteq B$ memberikan
$\dim_\Q B = \dim_\Q A \cdot \dim_A B$: sehingga faktor kirinya \omterm{def:b1:arith:divides}{membagi}.
Maka sublapangan $M$ berdimensi $\Q$ sebesar $1$, $2$ atau $4$:
dengan dimensi $1$ yaitu $\Q$ sendiri, dan dimensi $4$ yaitu $M$.

\textbf{23.} Misalkan $F \subseteq M$ sublapangan dengan $\dim_\Q F =
2$, dan $w \in F \setminus \Q$: maka $(1, w)$ \omterm{def:b1:vspaces:free}{bebas}, jadi sebuah \omterm{def:b1:vspaces:free}{basis}
$F$. Menurut pertanyaan 13 ($d = 2$), $w^2 = \alpha + \beta w$ bagi
rasional $\alpha, \beta$. Dengan menetapkan $v = w - \beta/2 \in F
\setminus \Q$:
\[
v^2 = w^2 - \beta w + \frac{\beta^2}4 = \alpha + \frac{\beta^2}4
\;\in\; \Q,
\]
dan $F = \operatorname{Vect}(1, v)$. Kini tulislah $v = a + b\sqrt2 +
c\sqrt3 + d\sqrt6$ lalu jabarkanlah:
\[
v^2 = (a^2 + 2b^2 + 3c^2 + 6d^2)
+ (2ab + 6cd)\sqrt2 + (2ac + 4bd)\sqrt3 + (2ad + 2bc)\sqrt6 .
\]
Ketiga \omterm{prop:b1:vspaces:coordinates}{koordinat} irasionalnya lenyap: $ab + 3cd = ac + 2bd =
ad + bc = 0$. Jika $a \neq 0$: maka $b = -3cd/a$, dan syarat
keduanya menjadi $c(a^2 - 6d^2) = 0$; dan karena $a^2 = 6d^2$ dengan $d
\neq 0$ akan membuat $\sqrt6 = \abs{a/d}$ rasional, maka $c = 0$
atau $d = 0$, dan pada kedua kasusnya syarat sisanya memaksa $b =
c = d = 0$, yaitu $v \in \Q$: yang dikecualikan. Jadi $a = 0$, dan
syaratnya terbaca $3cd = 2bd = bc = 0$: sehingga paling banyak satu di antara $b, c, d$ yang
taknol. Maka $v$ merupakan kelipatan rasional atas $\sqrt2$, $\sqrt3$
atau $\sqrt6$, dan
\[
F \in \bigl\{\Q(\sqrt2),\ \Q(\sqrt3),\ \Q(\sqrt6)\bigr\} ,
\]
yang masing-masingnya memang sublapangan berdimensi $2$ (dengan kestabilan seperti
pada pertanyaan 6, dan inversnya menurut pertanyaan 12): jadi tepat tiga
sublapangan kuadratik.

\textbf{24.} $(1 + \sqrt2 + \sqrt3)(1 + \sqrt2 - \sqrt3) = (1 +
\sqrt2)^2 - 3 = 2\sqrt2$, sehingga
\[
\frac1{1 + \sqrt2 + \sqrt3}
= \frac{1 + \sqrt2 - \sqrt3}{2\sqrt2}
= \frac{\sqrt2 + 2 - \sqrt6}{4}
= \frac12 + \frac14\sqrt2 + 0\cdot\sqrt3 - \frac14\sqrt6 .
\]
(Periksa: $(1 + \sqrt2 + \sqrt3)(2 + \sqrt2 - \sqrt6) = 4$ setelah
dijabarkan.)

\textbf{25.} (i) Dimensi $\Q$ melekatkan sebuah \emph{bilangan bulat} pada
setiap sublapangan $\R$, dan \omterm{pb:b1:findim:1}{hukum menara} membuat bilangan itu
mengalikan diri sepanjang pemuatannya: sehingga dimensinya berperilaku bagai
invarian aritmetika, dan kendala \omterm{def:b1:arith:divides}{keterbagiannya} menjadi bukti
kemustahilan. (ii) \omterm{def:b1:findim:def}{Dimensi hingga} memaksa setiap unsurnya memenuhi
persamaan \omterm{def:b1:poly:def}{polinomial} rasional taknol yang berderajat paling tinggi sebesar
dimensinya: jadi kehinggaan berarti kealjabaran. (iii) Adapun
kaidah dua-dari-tiga mengubah ``perkalian dengan $u$ memetakan sebuah \omterm{def:b1:vspaces:free}{basis}
ke \omterm{def:b1:vspaces:free}{keluarga bebas} berukuran maksimal'' menjadi kesurjektifan, yang menghasilkan
$u^{-1}$ tanpa rumus: sehingga inversnya lahir dari pencacahan. (iv) Adapun
\omterm{pb:b1:findim:1}{hukum menara} melarang \omterm{def:b1:structures:field}{lapangan} berdimensi $3$ di dalam yang
berdimensi $4$, dan itulah sebabnya $2^{1/3}$ tak dapat dicapai
dari $\sqrt2$ dan $\sqrt3$. Adapun teorema Bagian III adalah
\emph{\omterm{pb:b1:findim:1}{hukum menara} Dedekind}, yaitu langkah pembuka teori Galois,
yang dikembangkan pada jilid Tahun~3.

\end{solution}
